\documentclass[11pt,a4paper]{article}

% ── Packages ──────────────────────────────────────────────────────────────────
\usepackage[margin=2.5cm]{geometry}
\usepackage{amsmath,amssymb}
\usepackage{booktabs}
\usepackage{hyperref}
\usepackage{xcolor}
\usepackage{longtable}
\usepackage{array}
\usepackage{graphicx}
\usepackage{fancyhdr}
\usepackage{titlesec}
\usepackage{enumitem}

% ── Colour definitions ────────────────────────────────────────────────────────
\definecolor{apexblue}{RGB}{0,51,102}
\definecolor{apexgray}{RGB}{100,100,100}
\definecolor{apexred}{RGB}{180,30,30}

% ── Section formatting ────────────────────────────────────────────────────────
\titleformat{\section}{\large\bfseries\color{apexblue}}{}{0em}{\thesection\quad}
\titleformat{\subsection}{\normalsize\bfseries\color{apexblue}}{}{0em}{\thesubsection\quad}

% ── Header / footer ───────────────────────────────────────────────────────────
\pagestyle{fancy}
\fancyhf{}
\rhead{\textcolor{apexgray}{\small RESTRICTED --- Model Risk Sensitive}}
\lhead{\textcolor{apexgray}{\small APEX-MDL-0013 \textbar{} ColVA MDD v1.0}}
\rfoot{\textcolor{apexgray}{\small \thepage}}
\lfoot{\textcolor{apexgray}{\small Apex Global Bank \textbar{} Quantitative Research Division}}
\renewcommand{\headrulewidth}{0.4pt}
\renewcommand{\footrulewidth}{0.4pt}

% ── Hyperref setup ────────────────────────────────────────────────────────────
\hypersetup{
  colorlinks=true,
  linkcolor=apexblue,
  urlcolor=apexblue,
  pdftitle={Apex Global Bank -- ColVA Model Development Document},
  pdfauthor={XVA Quant Team}
}

% ══════════════════════════════════════════════════════════════════════════════
\begin{document}
% ══════════════════════════════════════════════════════════════════════════════

\begin{titlepage}
  \begin{center}
    \vspace*{2cm}
    {\color{apexblue}\rule{\linewidth}{2pt}}\\[0.5cm]
    {\Huge\bfseries\color{apexblue} Model Development Document}\\[0.4cm]
    {\LARGE\bfseries Collateral Valuation Adjustment (ColVA)}\\[0.3cm]
    {\Large Version 1.0}\\[0.5cm]
    {\color{apexblue}\rule{\linewidth}{2pt}}\\[1cm]
    {\large Apex Global Bank \textbar{} Quantitative Research Division}\\[0.3cm]
    {\large\textcolor{apexred}{Classification: RESTRICTED --- Model Risk Sensitive}}\\[0.3cm]
    {\large SR 11-7 Section Reference: Model Development}\\[2cm]
    \begin{tabular}{ll}
      \textbf{Model ID:}   & APEX-MDL-0013 \\[4pt]
      \textbf{Date:}       & March 2026 \\[4pt]
      \textbf{Status:}     & In Validation \\[4pt]
      \textbf{Risk Tier:}  & Tier 2 \\
    \end{tabular}
  \end{center}
\end{titlepage}

% ── Table of Contents ─────────────────────────────────────────────────────────
\tableofcontents
\newpage

% ══════════════════════════════════════════════════════════════════════════════
\section{Document Metadata}
% ══════════════════════════════════════════════════════════════════════════════

\begin{table}[h!]
\centering
\caption{Document Control Information}
\begin{tabular}{@{}ll@{}}
\toprule
\textbf{Field} & \textbf{Value} \\
\midrule
Model Name            & Collateral Valuation Adjustment (ColVA) \\
Model ID              & APEX-MDL-0013 \\
Version               & 1.0 \\
Status                & In Validation \\
Risk Tier             & Tier 2 \\
Regulatory Framework  & IFRS 13 Fair Value; ISDA Master Agreement / CSA \\
Date of Development   & March 2026 \\
Business Owner        & Head of Global Markets Trading (James Okafor) \\
Model Developer       & Quantitative Research --- XVA Team \\
MRO Review            & Pending --- Dr.\ Rebecca Chen \\
CRO Approval          & Pending --- Dr.\ Priya Nair \\
\bottomrule
\end{tabular}
\end{table}

\subsection{Version History}

\begin{table}[h!]
\centering
\caption{Document Version History}
\begin{tabular}{@{}lll@{}}
\toprule
\textbf{Version} & \textbf{Date} & \textbf{Change Summary} \\
\midrule
0.1 & Jan 2026 & Initial framework --- CTD currency optionality only \\
0.8 & Feb 2026 & Added rehypothecation value; multi-currency CSA \\
1.0 & Mar 2026 & Production candidate; submitted for validation \\
\bottomrule
\end{tabular}
\end{table}

% ══════════════════════════════════════════════════════════════════════════════
\section{Executive Summary}
% ══════════════════════════════════════════════════════════════════════════════

\subsection{Model Purpose}

ColVA (Collateral Valuation Adjustment) captures the present value of optionality embedded in
Credit Support Annexes (CSAs) governing collateral arrangements for OTC derivatives. The primary
source of ColVA is the \textbf{cheapest-to-deliver (CTD) option}: when a CSA permits the
collateral posting party to deliver collateral in any of several eligible currencies, that party
will always choose the currency that is cheapest in terms of its own funding cost relative to
the collateral rate paid. This embedded option has quantifiable economic value.

ColVA is classified as \textbf{Tier 2} within Apex's model inventory. While economically real,
it is smaller in magnitude and less data-intensive than the Tier~1 XVA models (CVA, FVA, MVA).
Nevertheless, for transactions where ColVA is material --- large cross-currency interest rate
swaps with multi-currency CSAs --- it must be priced into client derivatives quotes and
reflected in daily P\&L.

\subsection{Components of ColVA}

ColVA has three primary components:

\begin{enumerate}[leftmargin=*]
  \item \textbf{CTD Currency Option (dominant)}: The value to the posting party of being able to switch the currency of posted collateral to whichever is cheapest at the time of posting. For a 10-year EUR/USD cross-currency swap with a multi-currency CSA, this option can be worth 3--8 basis points of notional.

  \item \textbf{Substitution Option}: The right to substitute previously posted collateral (e.g., replace USD cash with EUR cash). This is only valuable when the relative cheapness of currencies changes after initial posting.

  \item \textbf{Rehypothecation Value}: When a CSA permits the collateral receiver to re-use posted collateral (e.g., post it as margin to a third party), the receiver captures additional funding benefit.
\end{enumerate}

\subsection{Output Description}

\begin{itemize}[leftmargin=*]
  \item \textbf{ColVA per new trade}: Priced into client derivatives quotes for transactions with multi-currency CSAs.
  \item \textbf{Portfolio ColVA}: Daily P\&L impact, reported in income statement under XVA line.
  \item \textbf{CTD option decomposition}: By currency pair, to inform collateral optimisation decisions.
  \item \textbf{ColVA by counterparty}: For counterparties with multi-currency CSA agreements.
\end{itemize}

\subsection{Key Assumptions}

\begin{enumerate}[leftmargin=*]
  \item The posting party will always exercise the CTD option optimally (rational collateral manager assumption).
  \item Currency basis swap spreads are the correct proxy for the cost differential between posting in different currencies.
  \item The collateral rate paid on cash collateral is the overnight index rate (SOFR, ESTR, SONIA) in the respective currency.
  \item Substitution rights are exercised only when the CTD currency changes (no partial exercise).
  \item The collateral pool is not subject to haircuts for the base calculation (haircut variant available separately).
\end{enumerate}

\subsection{Known Limitations}

\begin{itemize}[leftmargin=*]
  \item The model assumes the CSA counterparty does not strategically withhold consent to substitution --- in practice, substitution may be delayed or refused.
  \item Rehypothecation value is estimated using a simplified funding benefit; actual benefit depends on Apex's real-time collateral deployment strategy.
  \item The model does not capture operational costs of currency switching (FX transaction costs, settlement risk).
  \item Correlation between the portfolio's mark-to-market value and currency basis spreads is modelled via historical correlation but not stress-tested independently.
\end{itemize}

% ══════════════════════════════════════════════════════════════════════════════
\section{Business Context and Purpose}
% ══════════════════════════════════════════════════════════════════════════════

\subsection{The CSA Optionality Problem}

A Credit Support Annex specifies the legal terms under which collateral is posted and received
between counterparties. Many CSAs negotiated between Apex and its dealers specify that variation
margin may be posted in any of several eligible currencies --- typically USD, EUR, GBP, JPY, and
CHF. The posting party always posts the currency that costs it the least.

The cost of posting collateral in currency $X$ is:
\begin{equation}
\text{Posting cost}(X) = \text{Funding cost}(X) - \text{Collateral rate received}(X)
\tag{1}
\end{equation}

If Apex can fund USD at SOFR + 55\,bp but only receives SOFR on posted USD collateral, while it
can fund EUR at ESTR + 40\,bp and receives ESTR on EUR collateral, the cost differential is:
\begin{align*}
\text{USD cost} &= 55\,\mathrm{bp} \\
\text{EUR cost} &= 40\,\mathrm{bp} \quad\Rightarrow\quad \text{EUR is CTD by 15\,bp}
\end{align*}

Apex will post EUR, saving 15\,bp per annum on the collateral notional. Over the life of a
10-year swap, this saving --- and the option to switch currencies if the differential changes ---
is the ColVA.

\subsection{Market Context}

Multi-currency CSAs became the industry standard in the 2010s as banks sought to optimise their
collateral pools following the increased margin demands of the post-GFC regulatory environment.
ISDA's 2016 Credit Support Annex for Variation Margin (VM-CSA) standardised the framework, but
left currency eligibility as a negotiated term.

The theoretical framework for CTD option pricing was established by Fujii, Shimada, and
Takahashi (2010) and subsequently refined by Piterbarg (2010) in the context of derivatives
valuation under CSA.

\subsection{Model Users}

\begin{itemize}[leftmargin=*]
  \item \textbf{Front Office Pricing}: ColVA charge (or benefit, for the receiving party) included in the price of derivatives with multi-currency CSAs.
  \item \textbf{Collateral Management}: CTD decomposition informs which currency to post on a given day.
  \item \textbf{Finance / P\&L}: Daily ColVA P\&L recognised in income statement under XVA line.
\end{itemize}

% ══════════════════════════════════════════════════════════════════════════════
\section{Theoretical Foundation}
% ══════════════════════════════════════════════════════════════════════════════

\subsection{The CTD Option Value}

The value of the CTD option for a single time step $[t, t+\Delta t]$ in the risk-neutral
measure is:

\begin{equation}
\mathrm{CTD}(t) = \mathbb{E}\!\left[\max\!\left(c_1(t),\ldots,c_n(t)\right) - c_{\mathrm{base}}(t)\right]
  \cdot \mathbb{E}\!\left[C(t)\right] \cdot \Delta t \cdot D(0,t)
\tag{2}
\end{equation}

where $c_i(t)$ is the net funding benefit of posting in currency $i$ at time $t$,
$c_{\mathrm{base}}(t)$ is the funding benefit of posting in the base currency (USD), and
$C(t)$ is the expected collateral balance (expected absolute exposure under CSA terms).

The portfolio ColVA integrates over the life of all transactions:

\begin{equation}
\mathrm{ColVA} = \sum_i D(0,t_i)\cdot\mathbb{E}\!\left[\max\!\left(c_1(t_i),\ldots,c_n(t_i)\right) - c_{\mathrm{USD}}(t_i)\right]
\cdot \bigl|V(t_i)\bigr| \cdot \Delta t_i
\tag{3}
\end{equation}

where $|V(t)|$ is the absolute portfolio value (collateral balance under full collateralisation).

\subsection{Currency Basis Swap Spread as CTD Proxy}

The cost differential between posting in two currencies is proxied by the
\textbf{cross-currency basis swap spread}. For EUR vs.\ USD:

\begin{equation}
\Delta c_{\mathrm{EUR/USD}}(t) = s_{f,\mathrm{USD}}(t) - s_{f,\mathrm{EUR}}(t) - b_{\mathrm{EUR/USD}}(t)
\tag{4}
\end{equation}

where $b_{\mathrm{EUR/USD}}$ is the cross-currency basis (EUR/USD xccy basis; negative means
EUR is more expensive to source, reducing the CTD benefit of posting EUR).

Current market conditions (March 2026):

\begin{table}[h!]
\centering
\caption{Cross-Currency Basis Spreads (March 2026)}
\begin{tabular}{@{}ll@{}}
\toprule
\textbf{Currency Pair} & \textbf{xccy Basis} \\
\midrule
EUR/USD & $-15\,\mathrm{bp}$ \\
GBP/USD & $-8\,\mathrm{bp}$  \\
JPY/USD & $-45\,\mathrm{bp}$ \\
\bottomrule
\end{tabular}
\end{table}

\subsection{Stochastic Model for Basis Spreads}

Cross-currency basis spreads are mean-reverting with stochastic volatility. We model them using
an Ornstein-Uhlenbeck (OU) process:

\begin{equation}
d b(t) = \kappa\!\left(\theta - b(t)\right) dt + \sigma\, dW(t)
\tag{5}
\end{equation}

Parameters calibrated to 5-year historical data:

\begin{table}[h!]
\centering
\caption{OU Model Calibrated Parameters}
\begin{tabular}{@{}llll@{}}
\toprule
\textbf{Currency Pair} & \textbf{$\kappa$} & \textbf{$\theta$ (long-run mean)} & \textbf{$\sigma$ (bp$/\!\sqrt{\mathrm{yr}}$)} \\
\midrule
EUR/USD & 0.3 & $-12\,\mathrm{bp}$ & 8  \\
GBP/USD & 0.4 & $-7\,\mathrm{bp}$  & 6  \\
JPY/USD & 0.2 & $-38\,\mathrm{bp}$ & 12 \\
\bottomrule
\end{tabular}
\end{table}

The OU model is jointly simulated with the interest rate model (Hull-White, shared with CVA
simulation infrastructure) to capture the correlation between interest rate movements and
currency basis dynamics.

\subsection{Rehypothecation Value}

When Apex receives collateral under a CSA that permits rehypothecation, Apex can re-use that
collateral. The value of this right is:

\begin{equation}
V_{\mathrm{rehypo}} = \sum_i D(0,t_i)\cdot\mathbb{E}\!\left[C_{\mathrm{rcvd}}(t_i)\right]
\cdot \bigl(s_f - r_{\mathrm{OIS}}\bigr)\cdot\Delta t_i
\tag{6}
\end{equation}

where $s_f - r_{\mathrm{OIS}}$ is Apex's benefit from using received collateral rather than
borrowing in the unsecured market. At current funding spreads (55\,bp), the rehypothecation
value is modest but non-zero.

Rehypothecation is only relevant when Apex is the net receiver of collateral (i.e., when the
portfolio is in-the-money for Apex). For out-of-the-money positions, Apex is the poster, not
the receiver.

\subsection{Literature References}

\begin{itemize}[leftmargin=*]
  \item Piterbarg, V.\ (2010). ``Funding Beyond Discounting: Collateral Agreements and Derivatives Pricing.'' \textit{Risk Magazine}, February 2010.
  \item Fujii, M., Shimada, Y., Takahashi, A.\ (2010). ``A Note on Construction of Multiple Swap Curves with and without Collateral.'' \textit{CARF Working Paper}.
  \item Antonov, A., Bianchetti, M.\ (2011). ``Pricing Swaps and Options on Quadratic Gaussian Spreads.'' \textit{SSRN Working Paper}.
  \item Kenyon, C., Stamm, R.\ (2012). \textit{Discounting, LIBOR, CVA and Funding: Interest Rate and Credit Pricing}. Palgrave Macmillan.
  \item ISDA (2016). \textit{2016 Credit Support Annex for Variation Margin (VM)}.
\end{itemize}

% ══════════════════════════════════════════════════════════════════════════════
\section{Data Requirements and Governance}
% ══════════════════════════════════════════════════════════════════════════════

\subsection{Input Data Sources}

\begin{table}[h!]
\centering
\caption{ColVA Model Input Data Sources}
\begin{tabular}{@{}lll@{}}
\toprule
\textbf{Data Item} & \textbf{Source} & \textbf{Frequency} \\
\midrule
Cross-currency basis swap spreads      & Bloomberg / Markit             & Daily \\
OIS rates (SOFR, ESTR, SONIA, TONAR)  & Bloomberg / Internal           & Daily \\
Apex funding curve (all currencies)    & Treasury / Bloomberg           & Daily \\
CSA terms (eligible currencies, haircuts) & ISDA MA database            & On amendment \\
Portfolio exposures / absolute MtM     & CVA model Monte Carlo engine   & Daily (shared run) \\
FX spot rates                          & Bloomberg                      & EOD snapshot \\
\bottomrule
\end{tabular}
\end{table}

\subsection{CSA Terms Database}

The ColVA model is critically dependent on the accuracy of the CSA terms database, which
records:
\begin{itemize}[leftmargin=*]
  \item Eligible currencies for VM posting (by counterparty).
  \item Haircuts by currency and asset class.
  \item Rehypothecation rights (yes/no, by counterparty).
  \item Substitution notice periods (same-day, T$+$1, T$+$2).
\end{itemize}

The CSA terms database is maintained by the Legal/ISDA documentation team. Any change to a
counterparty's CSA terms must be reflected in the database within one business day. The ColVA
model validates against the database on each daily run and flags missing or stale CSA records.

\subsection{Data Governance}

ColVA shares the CVA model's simulation infrastructure (same paths, same exposure profiles).
Cross-currency basis spreads are sourced from Bloomberg (primary) and Markit (secondary). Any
day where the spread difference between sources exceeds 5\,bp triggers a data quality alert and
defaults to the prior-day value with a staleness flag.

% ══════════════════════════════════════════════════════════════════════════════
\section{Methodology and Implementation}
% ══════════════════════════════════════════════════════════════════════════════

\subsection{Shared Simulation Infrastructure}

ColVA runs as an extension of the CVA simulation:
\begin{enumerate}[leftmargin=*]
  \item CVA simulation produces $\mathbb{E}[|V(t)|]$ (expected absolute exposure) at each time step --- this is the collateral balance.
  \item ColVA applies the CTD option value to this collateral profile.
  \item The CTD option value at each time step comes from the OU basis spread model (jointly simulated paths).
\end{enumerate}

The joint simulation ensures that the correlation between market movements (which drive
portfolio MtM) and currency basis dynamics (which drive CTD optionality) is captured.

\subsection{Multi-Currency CSA Classification}

Counterparties are classified into three tiers for ColVA purposes:

\begin{table}[h!]
\centering
\caption{CSA Classification Tiers}
\begin{tabular}{@{}llp{7cm}@{}}
\toprule
\textbf{Tier} & \textbf{CSA Type} & \textbf{Treatment} \\
\midrule
A & USD-only CSA (no currency optionality) & ColVA $= 0$ \\
B & Bilateral multi-currency CSA (2--3 eligible currencies) & Full CTD model \\
C & ISDA Standard multi-currency (5+ eligible currencies) & Full CTD model with all pairs \\
\bottomrule
\end{tabular}
\end{table}

Approximately 35\% of Apex's derivatives notional is with Tier~B or~C counterparties. These
counterparties generate essentially all of Apex's ColVA.

\subsection{Collateral Optimisation vs.\ Valuation}

ColVA is a \textbf{valuation} model, not a \textbf{collateral optimisation} model. It answers
the question: ``What is the fair value of the CTD option?'' It assumes optimal future exercise
but does not instruct the collateral management team which currency to post today.

The Collateral Management team (under Treasury) runs a separate optimisation model to determine
the daily posting currency. ColVA and the optimisation model share data inputs but are governed
separately.

\subsection{Dealing with Haircuts}

When a CSA specifies haircuts (e.g., non-cash collateral such as government bonds receives a 2\%
haircut), the effective collateral balance is:

\begin{equation}
C_{\mathrm{eff}} = \mathrm{MtM}_{\mathrm{exposure}} \cdot \frac{1}{1 - h}
\tag{7}
\end{equation}

where $h$ is the haircut fraction. For example, posting a bond with a 2\% haircut requires
posting 102\% of the cash equivalent exposure. In practice, cash collateral (with no haircut)
dominates Apex's VM posting and the haircut adjustment is small.

% ══════════════════════════════════════════════════════════════════════════════
\section{Model Testing and Backtesting}
% ══════════════════════════════════════════════════════════════════════════════

\subsection{CTD Option Sensitivity}

The ColVA CTD option is most sensitive to:

\begin{table}[h!]
\centering
\caption{ColVA Key Sensitivities}
\begin{tabular}{@{}lp{9cm}@{}}
\toprule
\textbf{Factor} & \textbf{Effect} \\
\midrule
Basis spread level & $\pm$10\,bp shift in EUR/USD xccy basis changes portfolio ColVA by $\approx$\$4M \\
Basis spread volatility & Higher volatility increases the option value (convexity effect) \\
Portfolio duration & Longer-dated trades have more CTD optionality \\
\bottomrule
\end{tabular}
\end{table}

Sensitivity analysis conducted quarterly and reported to MRO.

\subsection{Benchmark Comparison}

A simplified closed-form approximation for the CTD option value uses the Margrabe (1978)
exchange option formula:

\begin{equation}
\mathrm{CTD}_{\mathrm{approx}} = C_{\mathrm{avg}} \cdot \mathrm{Margrabe}(s_1, s_2, \sigma_{\mathrm{basis}}, T)
\tag{8}
\end{equation}

where $\sigma_{\mathrm{basis}}$ is the volatility of the basis spread differential. Agreement
threshold: within 15\% of full simulation result. (Wider tolerance than CVA given the simpler
nature of the instrument being modelled.)

\subsection{P\&L Attribution}

Daily ColVA P\&L is attributed to:
\begin{itemize}[leftmargin=*]
  \item Change in basis spreads (delta).
  \item Change in portfolio exposure (new trades, maturities, market moves).
  \item Passage of time (theta).
\end{itemize}

P\&L explain $> 75\%$ of total ColVA daily change. Lower threshold than CVA/FVA because basis
spread moves can be idiosyncratic and difficult to predict intraday.

% ══════════════════════════════════════════════════════════════════════════════
\section{Performance Metrics and Calibration Standards}
% ══════════════════════════════════════════════════════════════════════════════

\begin{table}[h!]
\centering
\caption{ColVA Performance Metrics and Thresholds}
\begin{tabular}{@{}lll@{}}
\toprule
\textbf{Metric} & \textbf{Acceptable Range} & \textbf{Action if Breached} \\
\midrule
ColVA P\&L Explain          & $>75\%$                  & Investigate within 3 business days \\
Benchmark deviation (Margrabe) & $<15\%$ relative      & Review within 30 days \\
Basis spread staleness       & No spread $>3$ days old  & Data quality alert \\
CSA database completeness    & $>98\%$ active counterparties & Legal/documentation escalation \\
\bottomrule
\end{tabular}
\end{table}

\subsection{Recalibration Schedule}

\begin{itemize}[leftmargin=*]
  \item \textbf{Basis spread OU parameters}: Annual recalibration using rolling 5-year data.
  \item \textbf{Correlation (portfolio MtM vs.\ basis)}: Annual recalibration.
  \item \textbf{Simulation paths (shared with CVA)}: No independent recalibration; reviewed annually.
\end{itemize}

% ══════════════════════════════════════════════════════════════════════════════
\section{Limitations, Assumptions, and Known Weaknesses}
% ══════════════════════════════════════════════════════════════════════════════

\subsection{Optimal Exercise Assumption}

ColVA assumes the posting party always exercises the CTD option optimally. In practice:
\begin{itemize}[leftmargin=*]
  \item Collateral operations teams may not switch currencies daily due to operational friction.
  \item Currency substitution requires counterparty consent and may face delays.
  \item FX transaction costs eat into the basis spread advantage for small moves.
\end{itemize}

\textbf{Compensating control}: The model applies an operational discount factor of 85\% to
the theoretical CTD value, estimated from historical analysis of Apex's actual posting
behaviour vs.\ the optimal theoretical posting. This discount is reviewed annually.

\subsection{Rehypothecation Risk}

The rehypothecation value assumes Apex can always deploy received collateral at its full
funding benefit. In a stress scenario:
\begin{itemize}[leftmargin=*]
  \item Received collateral may be legally encumbered and unavailable for rehypothecation.
  \item The collateral receiver's funding benefit may differ from Apex's unsecured spread.
\end{itemize}

\textbf{Compensating control}: Rehypothecation value (currently $\approx$\$8M of total ColVA)
is disclosed separately and conservatively estimated.

\subsection{CSA Amendment Risk}

If a counterparty amends its CSA to remove multi-currency optionality (e.g., mandating USD-only
posting), the ColVA for that counterparty drops to zero immediately. A large counterparty
switching to a USD-only CSA could reduce portfolio ColVA by \$5--15M.

\textbf{Compensating control}: Legal team notifies the XVA desk within one business day of any
CSA amendment. Day-1 recalculation triggered automatically.

\subsection{Basis Spread Jump Risk}

The OU model captures mean-reversion but not jumps in currency basis spreads. During the
COVID-19 stress event (March 2020), EUR/USD xccy basis widened by 45\,bp in a single day ---
well outside the model's diffusive dynamics. This would cause ColVA to be underestimated during
stress events.

\textbf{Compensating control}: Stressed ColVA is computed using basis spreads at their 95th
percentile of historical 1-month widening. Disclosed in quarterly risk report.

% ══════════════════════════════════════════════════════════════════════════════
\section{Compensating Controls}
% ══════════════════════════════════════════════════════════════════════════════

\begin{table}[h!]
\centering
\caption{ColVA Compensating Controls}
\begin{tabular}{@{}p{4.5cm}p{6cm}p{3.5cm}@{}}
\toprule
\textbf{Risk} & \textbf{Control} & \textbf{Owner} \\
\midrule
Suboptimal exercise & 85\% operational discount factor; annual review & XVA Quant Team \\[4pt]
Basis spread staleness & 3-day staleness alert; prior-day fallback & Market Data \\[4pt]
CSA database gaps & 98\% completeness threshold; legal escalation & Legal / Documentation \\[4pt]
Basis jump risk & Stressed ColVA at 95th percentile spread widening & XVA Quant Team \\[4pt]
Rehypothecation overstatement & Separate disclosure; conservative estimate & XVA Quant / Treasury \\[4pt]
Shared simulation failure & Fallback: prior-day ColVA + basis delta & XVA Technology \\
\bottomrule
\end{tabular}
\end{table}

\subsection{Model Reserve}

A model reserve of \textbf{\$6 million} is held against ColVA model uncertainty:

\begin{table}[h!]
\centering
\caption{ColVA Model Reserve Breakdown}
\begin{tabular}{@{}lr@{}}
\toprule
\textbf{Reserve Component} & \textbf{Amount (USD M)} \\
\midrule
Optimal exercise discount factor uncertainty  & 3 \\
Basis spread model mis-specification          & 2 \\
CSA database completeness                     & 1 \\
\midrule
\textbf{Total Reserve}                        & \textbf{6} \\
\bottomrule
\end{tabular}
\end{table}

The ColVA reserve is small relative to CVA (\$45M) and FVA (\$25M), reflecting the Tier~2
classification of the model.

% ══════════════════════════════════════════════════════════════════════════════
\section{Relationship to Other XVA Models}
% ══════════════════════════════════════════════════════════════════════════════

\begin{table}[h!]
\centering
\caption{ColVA Interactions with Other XVA Models}
\begin{tabular}{@{}lp{10cm}@{}}
\toprule
\textbf{XVA Model} & \textbf{Interaction with ColVA} \\
\midrule
\textbf{CVA} & Shares Monte Carlo simulation paths and EPE/ENE profiles. \\[4pt]
\textbf{FVA} & FVA and ColVA both depend on the funding spread. No double-counting: FVA applies to uncollateralised positions; ColVA applies to optionality within collateralised positions. \\[4pt]
\textbf{MVA} & MVA applies to initial margin (uncleared derivatives); ColVA applies to variation margin optionality. Separate models with no double-counting. \\[4pt]
\textbf{DVA} & No direct interaction; ColVA does not depend on Apex's own default probability. \\
\bottomrule
\end{tabular}
\end{table}

The XVA aggregation report (produced daily) shows CVA $+$ DVA $+$ FVA $+$ MVA $+$ ColVA as a
single XVA total, with individual components disclosed separately for risk management purposes.

% ══════════════════════════════════════════════════════════════════════════════
\section{Change Management}
% ══════════════════════════════════════════════════════════════════════════════

Material changes (new CSA currency types, change to OU model specification, integration with
new simulation infrastructure) require full MRO re-validation. Non-material changes
(calibration updates, new counterparty CSA loaded into database) proceed via expedited review
with 3-business-day notification to MRO.

\subsection{Material Change Triggers}

\begin{itemize}[leftmargin=*]
  \item Addition of a new eligible currency class to any Tier~B or Tier~C CSA.
  \item Change to the OU basis spread model (e.g., addition of jump component).
  \item Change to the operational discount factor (currently 85\%).
  \item Replacement of the shared Monte Carlo simulation infrastructure.
  \item Change to rehypothecation accounting policy.
\end{itemize}

% ══════════════════════════════════════════════════════════════════════════════
\section{Approvals and Sign-Off}
% ══════════════════════════════════════════════════════════════════════════════

\begin{table}[h!]
\centering
\caption{Approval Status}
\begin{tabular}{@{}llll@{}}
\toprule
\textbf{Role} & \textbf{Name} & \textbf{Status} & \textbf{Date} \\
\midrule
Model Developer          & XVA Quant Team    & Submitted & March 2026 \\
Model Risk Officer       & Dr.\ Rebecca Chen & Pending   & --- \\
Chief Risk Officer       & Dr.\ Priya Nair   & Pending   & --- \\
Business Owner           & James Okafor      & Pending   & --- \\
\bottomrule
\end{tabular}
\end{table}

\vspace{1cm}
\noindent\textcolor{apexgray}{\footnotesize
Document Classification: RESTRICTED --- Model Risk Sensitive\\
Apex Global Bank \textbar{} Quantitative Research Division \textbar{} XVA Team\\
Next scheduled review: March 2027 (or upon material model change)
}

\end{document}
